% !TeX root = ../../Book1.tex
\section{Morrey 理论}
\label{b1:sec:2403}

当梯度的可积指数超过空间维数时，局部平均振荡不仅趋零，而且在逐次缩小半径后可以求和。这使 Sobolev 函数获得连续代表，并给出两点之间的定量控制。以下先讨论 \(d\ge1\)、\(d<p<\infty\)，令
\[
 \alpha=1-\frac dp\in(0,1).
\]

为把本节放回整章的位置，\cref{b1:tab:first-order-sobolev-regimes}列出 \(d\ge2\) 时一阶理论的三个指数区间。它只给出典型目标；区域、边界条件与端点的精确假设仍以各正式定理为准。
\begin{table}[htbp]
  \centering
  \caption[一阶 Sobolev 的三个区间]{一阶 Sobolev 的三个区间。梯度尺度从积分增益，经临界端点，过渡到 Hölder 连续性。}
  \label{b1:tab:first-order-sobolev-regimes}
  \begin{tabularx}{\textwidth}{@{}p{2.7cm}p{3.6cm}X@{}}
    \toprule
    指数范围 & 临界尺度 & 典型结论 \\
    \midrule
    \(1\le p<d\) & \(p^*=dp/(d-p)\) & 嵌入到 \(L^{p^*}\)，紧性只在严格低于临界指数时出现 \\
    \(p=d\) & \(p^*=\infty\) 的形式端点 & 一般无 \(L^\infty\) 控制；有界域上有全部有限 \(L^q\) 控制，零边界时还有指数可积性 \\
    \(p>d\) & \(\alpha=1-d/p>0\) & 存在 \(C^{0,\alpha}\) 代表，并有 Morrey 型双点估计 \\
    \bottomrule
  \end{tabularx}
\end{table}

以下以实值函数为主，复值或有限维向量值函数可按分量处理。记 \(\|\nabla u\|_{L^p}\) 为弱梯度的欧氏长度的 \(L^p\) 范数；它与本书按偏导分量计算的梯度范数只相差依赖维数及 \(p\) 的等价常数。

Hajłasz 的《Sobolev spaces on metric-measure spaces》\cite[Section 9]{Hajlasz2003SobolevMetric}用球均值的望远镜求和把平均振荡变成双点估计。下面保留这一组织，但利用欧氏体积与 \(p>d\)，直接让球均值在每个点收敛，不先把讨论限制在 Lebesgue 点上。
% 实读 Hajlasz2003SobolevMetric：作者正式公开PDF，§9 Theorem 9.4及完整证明，
% 印刷204--205页；另读Theorem 9.7(3)的Hölder结论陈述，印刷206页。
% 这里改编逐半径球均值比较；一般度量空间的先修与9.7完整证明不作为本节依赖。
% 欧氏球Poincare由本章24.1提供；每一点的极限、任意半径一致性及代表辨析另行展开。

\subsection{\texorpdfstring{\(p>d\)}{p>d}}
\label{b1:subsec:240301}

对球 \(B=B(x,r)\)，记 \(u_B=m_d(B)^{-1}\int_Bu\)。球上的 Poincaré 不等式与 Hölder 不等式给出
\begin{equation}\label{b1:eq:morrey-average-oscillation}
 \begin{split}
 \frac1{m_d(B)}\int_B|u-u_B|
 &\le m_d(B)^{-1/p}\|u-u_B\|_{L^p(B)}\\
 &\le C_{d,p}r^{1-d/p}\|\nabla u\|_{L^p(B)}.
 \end{split}
\end{equation}
这里使用的是\cref{b1:thm:ball-poincare-lp}，无需先知道 \(u\) 连续。幂次 \(1-d/p\) 是导数贡献的一个长度因子，与平均化引入的 \(r^{-d/p}\) 相乘所得。当它为正时，各个尺度的误差形成收敛的几何级数。

\begin{lemma}\label{b1:lem:morrey-ball-average-limit}
若 \(u\in W^{1,p}_{\mathrm{loc}}(\Omega)\)，其中 \(\Omega\subset\mathbb R^d\) 为开集、\(d<p<\infty\)，则对每个 \(x\in\Omega\)，有限极限
\[
 u^*(x)=\lim_{r\downarrow0}u_{B(x,r)}
\]
都存在，且 \(u^*=u\) 几乎处处。对每个满足 \(\overline{B(x,R)}\Subset\Omega\) 的球，还有
\begin{equation}\label{b1:eq:morrey-average-error}
 |u^*(x)-u_{B(x,R)}|
 \le C_{d,p}R^\alpha\|\nabla u\|_{L^p(B(x,R))}.
\end{equation}
\end{lemma}
\begin{proof}
固定 \(x,R\)，置 \(r_j=2^{-j}R\) 和 \(B_j=B(x,r_j)\)。因为两相邻球的体积比为 \(2^d\)，
\[
 \begin{split}
 |u_{B_{j+1}}-u_{B_j}|
 &\le\frac1{m_d(B_{j+1})}\int_{B_j}|u-u_{B_j}|\\
 &\le C_{d,p}r_j^\alpha\|\nabla u\|_{L^p(B_j)}
 \le C_{d,p}2^{-j\alpha}R^\alpha
                         \|\nabla u\|_{L^p(B_0)}.
 \end{split}
\]
右端可求和，所以均值序列为 Cauchy 列，存在有限极限；从 \(j=0\) 起求和就得到\eqref{b1:eq:morrey-average-error}。

还要排除极限对选定半径列的依赖。若 \(r_{j+1}<r\le r_j\)，则 \(B(x,r)\subset B_j\)，且体积比仍不超过 \(2^d\)。同一估计给出
\[
 |u_{B(x,r)}-u_{B_j}|
 \le C_{d,p}r_j^\alpha\|\nabla u\|_{L^p(B_0)}
 \longrightarrow0.
\]
因此所有趋零半径的均值都趋于同一极限，它与最初的 \(R\) 无关。以上论证对每个固定点成立，没有删除例外点。最后由\cref{b1:thm:lebesgue-point}，在几乎每个点，该极限等于原函数值，故确实得到 \(u\) 的一个代表。
\end{proof}

当 \(p=d\) 时，上面用整个固定球上的梯度范数估计的几何级数失去衰减。这个观察只说明当前证明不能给出正的 Hölder 指数。\secref{b1:sec:2405}将讨论临界情形中的其他控制。

\subsection{全空间 Morrey 不等式}
\label{b1:subsec:240302}

\begin{definition}\label{b1:def:holder-space}
给定非空集合 \(A\subset\mathbb R^d\) 及 \(0<\gamma\le1\)，若函数 \(v:A\rightarrow\mathbb R\) 满足
\[
 [v]_{0,\gamma;A}
 \coloneqq\sup_{\substack{x,y\in A\\x\ne y}}
       \frac{|v(x)-v(y)|}{|x-y|^\gamma}<\infty,
\]
则称它在 \(A\) 上为 \(\gamma\) 阶\term[holder-lianxu]{Hölder 连续}[Hölder continuous]。若还满足 \(\sup_A|v|<\infty\)，记 \(v\in C^{0,\gamma}(A)\)，并规定
\[
 \|v\|_{C^{0,\gamma}(A)}
 \coloneqq\sup_A|v|+[v]_{0,\gamma;A}.
\]
集合只有一个点时，半范数中的上确界约定为零。
\end{definition}
\symidx{\(C^{0,\gamma}(A)\)：有界Hölder连续函数空间}
\symidx{\([v]_{0,\gamma;A}\)：Hölder半范数}

指数为 \(1\) 时就是 Lipschitz 条件。对紧集，有限半范数和任意一个有限点值已保证有界；在全空间上则不能省去零阶条件，例如 \(v(x)=x_1\) 虽然 Lipschitz，却不属于这里的 \(C^{0,1}(\mathbb R^d)\)。

\begin{theorem}[Morrey]\label{b1:thm:morrey-whole-space}
设 \(d<p<\infty\) 且 \(u\in W^{1,p}(\mathbb R^d)\)。球平均极限 \(u^*\) 是唯一的连续代表，属于 \(C^{0,\alpha}(\mathbb R^d)\)，其中 \(\alpha=1-d/p\)，并满足
\begin{equation}\label{b1:eq:morrey-whole-space}
 [u^*]_{0,\alpha;\mathbb R^d}
 \le C_{d,p}\|\nabla u\|_{L^p(\mathbb R^d)},
 \quad
 \|u^*\|_{C^{0,\alpha}(\mathbb R^d)}
 \le C_{d,p}\|u\|_{W^{1,p}(\mathbb R^d)}.
\end{equation}
\end{theorem}
\begin{proof}
给定 \(x\ne y\)，令 \(r=|x-y|\)，并取
\[
 B_x=B(x,r),\quad B_y=B(y,r),\quad B_*=B(x,3r).
\]
两个小球都包含于 \(B_*\)，体积比只依赖于 \(d\)。由\eqref{b1:eq:morrey-average-oscillation}，
\[
 |u_{B_x}-u_{B_*}|+|u_{B_y}-u_{B_*}|
 \le C_{d,p}r^\alpha\|\nabla u\|_{L^p(B_*)}.
\]
再在 \(x,y\) 两点使用\eqref{b1:eq:morrey-average-error}，得到
\begin{equation}\label{b1:eq:morrey-local-pair}
 |u^*(x)-u^*(y)|
 \le C_{d,p}|x-y|^\alpha\cdot
                 \|\nabla u\|_{L^p(B(x,3|x-y|))}.
\end{equation}
放大积分域到整个空间，即得 Hölder 半范数界，也证明 \(u^*\) 处处连续。

半范数不能单独控制常数，因此还须估计一个平均值。对任意 \(x\)，用半径为 \(1\) 的球，
\[
 \begin{split}
 |u^*(x)|
 &\le|u^*(x)-u_{B(x,1)}|+|u_{B(x,1)}|\\
 &\le C_{d,p}\|\nabla u\|_{L^p(\mathbb R^d)}
       +\omega_d^{-1/p}\|u\|_{L^p(\mathbb R^d)}.
 \end{split}
\]
于是得到全范数估计。若另一个连续函数与 \(u\) 几乎处处相等，它与 \(u^*\) 的差连续且几乎处处为零；一旦在某点非零，就会在一个正测度小球内非零，矛盾。因此连续代表唯一。
\end{proof}

上述\term[morrey-budengshi]{Morrey 不等式}[Morrey inequality]把弱梯度的积分控制转成点值的变化控制。证明中的每一个平均值只依赖等价类，连续代表不是通过任意指定某个零测集上的值获得的。

\subsection{局部与区域版本}
\label{b1:subsec:240303}

全空间估计也给出内部估计。这里必须留意参与均值比较的球是否仍在原开集内，而不能直接对靠近边界的点使用越出开集的球。

\begin{corollary}\label{b1:cor:morrey-domain}
设 \(d<p<\infty\)、\(\alpha=1-d/p\)。
若 \(u\in W^{1,p}_{\mathrm{loc}}(\Omega)\)，则其球平均代表局部 Hölder 连续。具体地，若 \(\overline{B(a,8R)}\Subset\Omega\)，则
\begin{equation}\label{b1:eq:morrey-interior-ball}
 \begin{split}
 [u^*]_{0,\alpha;\overline{B(a,R)}}
 &\le C_{d,p}\|\nabla u\|_{L^p(B(a,8R))},\\
 \sup_{\overline{B(a,R)}}|u^*|
 &\le C_{d,p}\left(
        R^{-d/p}\|u\|_{L^p(B(a,8R))}
        +R^\alpha\|\nabla u\|_{L^p(B(a,8R))}\right).
 \end{split}
\end{equation}
若 \(\overline U\Subset V\subset\Omega\)，其中 \(U,V\) 为开集，且 \(u\in W^{1,p}(V)\)，则
\[
 \|u^*\|_{C^{0,\alpha}(\overline U)}
 \le C_{U,V,d,p}\|u\|_{W^{1,p}(V)}.
\]
若 \(\Omega\) 为有界 Lipschitz 区域，则 \(u\in W^{1,p}(\Omega)\) 有唯一的连续延伸到 \(\overline\Omega\) 的代表，且
\begin{equation}\label{b1:eq:morrey-lipschitz-domain}
 \|u^*\|_{C^{0,\alpha}(\overline\Omega)}
 \le C_{\Omega,p}\|u\|_{W^{1,p}(\Omega)}.
\end{equation}
\end{corollary}
\begin{proof}
对 \(x,y\in\overline{B(a,R)}\)，有 \(|x-y|\le2R\)，故
\(B(x,3|x-y|)\subset B(a,8R)\)。主定理的双点比较只在这些内部球中进行，所以\eqref{b1:eq:morrey-local-pair}仍成立，给出第一个估计。对 \(x\) 使用 \(B(x,R)\) 的平均值，并以 Hölder 不等式控制该平均值，就得到第二个估计。这同时证明代表的局部 Hölder 连续性。

给定 \(\overline U\Subset V\)，取固定 \(\chi\in C_c^\infty(V)\)，使它在 \(\overline U\) 的一个邻域内等于 \(1\)。将 \(\chi u\) 补零到整个空间。由弱乘积公式及\cref{b1:lem:compact-sobolev-zero-extension}，所得 \(w\) 满足
\[
 \|w\|_{W^{1,p}(\mathbb R^d)}
 \le C_{\chi,p}\|u\|_{W^{1,p}(V)}.
\]
在 \(\overline U\) 附近的充分小球上，\(u\) 与 \(w\) 的平均值相同，故两者的精确代表也相同。应用全空间估计即得第二项结论。

最后，若 \(\Omega\) 有界且为 Lipschitz 区域，用\cref{b1:thm:lipschitz-domain-extension}取得全空间延拓 \(Eu\)，再把其 Hölder 连续代表限制到 \(\overline\Omega\)。在内部，它与原来的球平均代表一致；边界值由内部连续趋近唯一确定。因此结果不依赖所选延拓算子，且延拓范数界给出\eqref{b1:eq:morrey-lipschitz-domain}。
\end{proof}

区域不必连通，但区域版本使用的是完整 \(W^{1,p}\) 范数，不能一律只保留 \(\|\nabla u\|_p\)。例如在两个分离的球上分别取常数 \(0\) 与 \(1\)，弱梯度为零，跨两个分量的 Hölder 半范数却非零。全空间延拓同时处理各分量之间的变化，相关常数包含区域的几何信息。

当 \(p=\infty\) 时，结论中的指数改为 \(1\)，但证明不通过把 \(p\) 代入有限指数公式来完成。\cref{b1:thm:w1infty-lipschitz}已经给出局部 Lipschitz 代表；在全空间中，把连接两点的线段分为有限个落在局部球内的小段并求和，就得到由全域弱梯度本性界控制的 Lipschitz 常数。其上确界由 \(\|u\|_\infty\) 控制。再使用截断或有界 Lipschitz 区域的延拓算子，分别得到局部和闭区域上的 \(C^{0,1}\) 估计。

指数 \(\alpha=1-d/p\) 在统一范数估计中不能提高。固定非恒定函数 \(\phi\in C_c^\infty(B(0,1))\)，并令
\[
 u_\varepsilon(x)=\varepsilon^\alpha\phi(x/\varepsilon),
 \quad 0<\varepsilon<1.
\]
换元直接给出
\[
 \|u_\varepsilon\|_p=\varepsilon\|\phi\|_p,\quad
 \|\nabla u_\varepsilon\|_p=\|\nabla\phi\|_p.
\]
这些函数在 \(W^{1,p}\) 中一致有界。然而，选取 \(\phi(a)\ne\phi(b)\)，对每个 \(\alpha<\beta\le1\)，
\[
 [u_\varepsilon]_{0,\beta;\mathbb R^d}
 \ge\varepsilon^{\alpha-\beta}
       \frac{|\phi(a)-\phi(b)|}{|a-b|^\beta}
 \longrightarrow\infty.
\]
所以不存在把所有 \(W^{1,p}\) 函数的 \(C^{0,\beta}\) 范数控制住的统一常数；同一例子也适用于包含这些小支集的固定球。每个 \(u_\varepsilon\) 本身仍然光滑，这里排除的是更强的连续嵌入估计，而不是说单个函数不能具有更好的正则性。

\subsection{精确代表}
\label{b1:subsec:240304}

\secref{b1:sec:1404}曾由球平均定义精确代表，并在极限不存在处另作赋值。现在，超临界可积性使这一例外情形在内部完全消失。

\begin{proposition}\label{b1:prop:morrey-precise-representative}
若 \(u\in W^{1,p}_{\mathrm{loc}}(\Omega)\)、\(d<p<\infty\)，则其精确代表 \(u^*\) 是唯一的连续代表，并且对每个 \(x\in\Omega\)，
\[
 \lim_{r\downarrow0}\frac1{m_d(B(x,r))}
       \int_{B(x,r)}|u(y)-u^*(x)|\,\mathrm dy=0.
\]
特别地，\(u^*\) 的每个内部点都是 Lebesgue 点。若 \(u_j\to u\) 于 \(W^{1,p}_{\mathrm{loc}}(\Omega)\)，则对每个相对紧开集 \(U\)，精确代表在 \(C^{0,\alpha}(\overline U)\) 中收敛；这里取 \(\overline U\Subset\Omega\)，且 \(\alpha=1-d/p\)。
\end{proposition}
\begin{proof}
有限极限已由\cref{b1:lem:morrey-ball-average-limit}在每个点构造，局部 Hölder 连续性由前一推论给出。连续代表的唯一性仍由非空开球具有正测度得到。对充分小的 \(r\)，
\[
 \begin{split}
 \frac1{m_d(B(x,r))}\int_{B(x,r)}|u-u^*(x)|
 &\le\frac1{m_d(B(x,r))}\int_{B(x,r)}|u-u_{B(x,r)}|
       +|u_{B(x,r)}-u^*(x)|\\
 &\le C_{d,p}r^\alpha
                 \|\nabla u\|_{L^p(B(x,r))}
 \longrightarrow0.
 \end{split}
\]
将积分中的 \(u\) 换成与它几乎处处相等的 \(u^*\)，即得每点的 Lebesgue 点性质。

最后取 \(\overline U\Subset V\) 且 \(\overline V\Subset\Omega\)。球均值的线性及每点极限说明
\((u_j-u)^*=u_j^*-u^*\)。对差函数使用前一推论，便有
\[
 \|u_j^*-u^*\|_{C^{0,\alpha}(\overline U)}
 \le C_{U,V,d,p}\|u_j-u\|_{W^{1,p}(V)}
 \longrightarrow0.
\]
\end{proof}

连续性属于这个被确定的代表，而不是等价类中的每个可测函数。即使只在一点改变 \(u^*\) 的值，也不改变任何 \(L^p\) 范数或弱导数，却可能破坏该点的连续性。另一方面，在有界 Lipschitz 区域上，\cref{b1:thm:lipschitz-lp-trace}说明这个闭域连续代表的边界限制恰是前章的迹；积分方式定义的边界值与通常点值在这里相遇。

本节得到的是函数值的 Hölder 控制，并不使弱梯度自动连续。把同一方法应用到哪些低阶导数、怎样组织高阶连续代表，是\secref{b1:sec:2404}高阶嵌入要处理的问题。
